\section{The upper comparison}\label{sec:upper}

We prove the upper bound in Theorem~\ref{thm:main} for a finite
spectral law
\(\mu=\sum_{a=1}^m\rho_a\delta_{\lambda_a}\), with every
\(\rho_a>0\). The spectral counts may be any deterministic sequence
whose proportions converge to these masses. The rotation identities
express the limiting energy through the total-overlap quantile. We
first identify an order in which this energy is monotone, then obtain
that order by minimizing an enriched pressure against a fixed trial.

\begin{lemma}[Energy and tail order]\label{up:energy}
For \(t\in[0,1]\) and \(p\in\cP\), define
\[
 F_t(x)=xR(tx),\qquad
 \mathcal E_t(p)=\frac12\int_0^1F_t'(D_p(r))\,dr.
\]
At a limit satisfying Proposition~\ref{rot:paths},
\(\mathcal E_t(p)\) is the expected interaction energy per spin
with coefficient \(\Lambda_N\), before multiplication by \(t\).
The function \(F_t\) is convex on \([0,1]\). In particular,
\begin{equation}\label{up:tail-order}
 \int_s^1p(u)\,du\ge\int_s^1q(u)\,du
 \quad(0\le s\le1)
 \quad\Longrightarrow\quad
 \mathcal E_t(p)\le\mathcal E_t(q).
\end{equation}
Moreover,
\begin{equation}\label{up:primitive}
 \int_0^t\mathcal E_u(q)\,du
 =G_t(q)=\frac12\int_0^1tR(tD_q(r))\,dr.
\end{equation}
Writing \(L=\max_a|\lambda_a|\), one has
\(\lvert\mathcal E_t(p)\rvert\le L/2\).
\end{lemma}

\begin{proof}
Let \(b_t\) and \(f_a\) be the resolvent functions for the spectrum
\(t\lambda_a\), as in Proposition~\ref{rot:paths}. That proposition
gives
\[
 \frac12\sum_a\lambda_a d_a
 =\frac12\int_0^1\sum_a\lambda_a f_a'(D_p(r))\,dr.
\]
For \(t>0\), the identity \(b_t(x)=t b(tx)\) implies
\[
 \sum_a\lambda_a f_a(x)
 =\frac{x b_t(x)-1}{t}=xR(tx)=F_t(x).
\]
The same identity holds by continuity at \(t=0\) and at \(x=0\).
This proves the energy formula.

For \(t>0\) and \(x>0\), set
\[
 M_j=\sum_a\frac{\rho_a}{(b_t(x)-t\lambda_a)^j}.
\]
The resolvent equation says \(M_1=x\), and differentiation yields
\[
 b_t'(x)=-\frac1{M_2},\qquad
 b_t''(x)=\frac{2M_3}{M_2^3},\qquad
 (x b_t(x))''=\frac{2(M_1M_3-M_2^2)}{M_2^3}\ge0.
\]
The inequality is Cauchy--Schwarz. Smooth extension at zero proves
convexity on the closed interval; at \(t=0\), \(F_0\) is linear.
Also \(f_a'\ge0\) and \(\sum_a f_a'=1\), so the displayed energy
is an average of numbers in \([-L/2,L/2]\).

For a nondecreasing quantile, its positive-part integral can be
recovered from upper tails:
\[
 \int_0^1(p(u)-r)_+\,du
 =\max_{0\le s\le1}
   \left\{\int_s^1p(u)\,du-(1-s)r\right\}.
\]
The hypothesis of \eqref{up:tail-order} therefore gives
\(D_p\le D_q\). Since \(F_t'\) is nondecreasing, the energy order
follows. Finally,
\[
 \frac{\partial}{\partial t}\bigl[tR(tx)\bigr]
 =R(tx)+txR'(tx)=F_t'(x).
\]
Integration in \(t\) proves \eqref{up:primitive}, including its
endpoint at zero.
\end{proof}

\begin{proposition}[Finite-spectrum upper bound]\label{up:bound}
Suppose the deterministic eigenvalues belong to a fixed finite set
\(\{\lambda_1,\ldots,\lambda_m\}\), and their proportions converge
to \(\rho_a>0\). Then
\[
 \limsup_{N\to\infty}\E P_N
 \le\inf_{p\in\cP}\{S(p)+G_1(p)\}.
\]
\end{proposition}

\begin{proof}
Fix a finite step trial \(q\) with values strictly below one. Denote
its interval weights by \(w_i>0\), its values by \(q_i\), and its
left endpoints by \(\zeta_i\), for \(0\le i<k\). By
Lemma~\ref{fld:domain}, \(S(q)<\infty\). Use the cascade on these
intervals to form the enriched, perturbed pressure
\(P_N^{\mathrm{enr}}(t,h,\mathbf u,\mathbf v)\) of
Section~\ref{sec:rotation}. Its linear field has covariance levels
\(h_i\), and the exponent includes the subtraction
\(-Nh_{k-1}/2\).

Suppose that the upper bound \(S(q)+G_1(q)\) fails along a
subsequence. Choose \(\delta>0\) small enough that, on this
subsequence, the objective
\begin{equation}\label{up:objective}
 \begin{aligned}
 \mathcal J_N(t,h,\mathbf u,\mathbf v)
 ={}&-\E P_N^{\mathrm{enr}}(t,h,\mathbf u,\mathbf v)
  -\frac12\int_0^1q(s)h(s)\,ds\\
 &+S(q)+G_t(q)+t\delta+\Pi_N(\mathbf u,\mathbf v).
 \end{aligned}
\end{equation}
has value at most \(-\eta\), for some fixed \(\eta>0\), at
\(t=1\), \(h=0\), and all perturbation parameters equal to
\(3/2\). Here
\[
 \Pi_N(\mathbf u,\mathbf v)
 =\sum_{j\le N}2^{-j}(u_j-3/2)^2
  +\sum_a(v_a-3/2)^2.
\]
This follows from the \(O(e_N)\) perturbation comparison. Minimize
\(\mathcal J_N\) on
\[
 0\le t\le1,\qquad
 0\le h_0\le\cdots\le h_{k-1}\le H,\qquad
 \mathbf u\in[1,2]^N,\quad\mathbf v\in[1,2]^m,
\]
where the fixed cap \(H\) will be chosen below. Gaussian and direct
comparison make the objective continuous, so a minimum exists.

At a negative minimum, increasing the field in nonnegative monotone
directions will force the limiting overlap quantile to have upper tails
at least as large as those of \(q\). Lemma~\ref{up:energy} will then
order their energies. A feasible decrease in \(t\) gives the
opposite energy inequality, with the strict margin \(\delta\).
To make these two variations available, we first exclude \(t=0\)
and the field cap. At \(t=0\), the unperturbed enriched pressure is exactly
\(\ell(h)\). The definition of \(S(q)\) and the uniform perturbation
bound give
\begin{equation}\label{up:zero-time}
 \mathcal J_N(0,h,\mathbf u,\mathbf v)\ge-Ce_N.
\end{equation}
In fact the minimizing values of \(t\) stay away from zero. For fixed
other parameters, Lemma~\ref{up:energy} and the spin norm bound give
\[
 \left|\partial_t\mathcal J_N\right|\le L+\delta.
\]
Comparing a minimum of value at most \(-\eta\) with
\eqref{up:zero-time}, while retaining its other parameters, yields
\(t_N\ge\eta/[2(L+\delta)]\) for all sufficiently large \(N\).

For the field cap, the recursion and \(\log\cosh z\le|z|\) give
\begin{equation}\label{up:cap-estimate}
 \ell(h)\le
 -\frac12\sum_iw_i h_i+\sqrt{\frac{2h_0}{\pi}}+C_q.
\end{equation}
Indeed a Gaussian increment of variance \(v\) at positive exponent
\(\zeta_i\) sends \(|z|\) to a function bounded by
\(|z|+\zeta_i v/2+(\log2)/\zeta_i\). Iterate over the positive
levels and take ordinary expectation at the root. The identity
\(\sum_i\zeta_i(h_i-h_{i-1})=h_{k-1}-\sum_iw_i h_i\)
then proves \eqref{up:cap-estimate}; when \(k=1\), the sum of
positive-level constants is empty.

The interaction changes the normalized log partition by at most
\(L/2\), so
\[
 \E P_N^{\mathrm{enr}}(t,h,\mathbf u,\mathbf v)
 \le\ell(h)+L/2+Ce_N.
\]
Put \(\epsilon=1-\max_iq_i>0\). Since \(S(q)\ge0\),
\(G_t(q)\ge-L/2\), and the penalty is nonnegative,
\eqref{up:cap-estimate} implies, on \(h_{k-1}=H\),
\[
 \mathcal J_N
 \ge\frac{\epsilon w_{k-1}}2H
       -\sqrt{\frac{2H}{\pi}}-C_q-L-Ce_N.
\]
Choose \(H\) large enough to make this positive for large \(N\).
Thus every minimum under consideration has \(h_{k-1}<H\).
The partition of the trial is fixed throughout this choice; no
uniformity in its number of intervals is required.

At these minima, Proposition~\ref{rot:gg} applies to the deterministic
minimizing perturbation parameters. Pass to a subsequence so that
\(t\) and the overlap arrays converge, and let \(p\)
be the limiting quantile of \(Q_{12}\). The resulting array satisfies
joint GG for the projected overlaps and the cascade level \(T\).
The cascade marking identity preserves the level probabilities
\(w_i\). Proposition~\ref{arr:structure} therefore represents
\(Q_{12}\) and \(T_{12}\) by nondecreasing quantiles on the same
unit interval. The quantile \(p\) need not be constant on the trial
intervals.

We now use the field coordinates of the minimum to order the upper
tails of \(p\) and \(q\). Let \(g\) be any nonnegative
nondecreasing step function on the fixed trial intervals, with values
\(g_i\), and define \(\gamma_g(i/k)=g_i\) on the cascade levels.
A small
increase of \(h\) in direction \(g\) is feasible, even if some
levels are zero or tied, because the upper cap is inactive. The
covariance derivative of the linear field is
\(NQ_{12}\gamma_g(T_{12})\). Its diagonal term is canceled by the
subtraction \(-Nh_{k-1}/2\), so minimality gives
\[
 0\le D_g\mathcal J_N
 =\frac12\E\langle Q_{12}\gamma_g(T_{12})\rangle
  -\frac12\int_0^1q(s)g(s)\,ds.
\]
By synchronized quantiles, the limit of the first expectation is
\(\int_0^1p(s)g(s)\,ds\). Hence
\begin{equation}\label{up:contact-order}
 \int_0^1p(s)g(s)\,ds\ge\int_0^1q(s)g(s)\,ds.
\end{equation}

Choose \(g\) to be each upper-tail indicator of the partition.
This proves the tail inequalities at its endpoints. The function
\(s\mapsto\int_s^1p\) is concave, because \(p\) is nondecreasing,
and therefore lies above its chords. The corresponding tail
integral of \(q\) is affine on each trial interval. The endpoint
inequalities consequently hold at every \(s\in[0,1]\).
Lemma~\ref{up:energy} now yields
\(\mathcal E_t(p)\le\mathcal E_t(q)\).

It remains to vary \(t\) at the same minimum.
A move to the left is feasible, including when \(t_N=1\), so
\(\partial_t\mathcal J_N\le0\). Proposition~\ref{rot:paths}
and Lemma~\ref{up:energy} identify the limit of this derivative as
\[
 -\mathcal E_t(p)+\mathcal E_t(q)+\delta\ge\delta>0.
\]
This contradiction proves the upper bound for the fixed step
trial \(q\).

Finally, given any \(p\in\cP\) with \(S(p)<\infty\), round
its values down to finite steps \(q_n\) such that
\(q_n\le p\), \(\max q_n<1\), and
\(\|q_n-p\|_\infty\to0\). Nonnegativity of every admissible field
gives \(S(q_n)\le S(p)\), while
\[
 \|D_{q_n}-D_p\|_\infty\le\|q_n-p\|_1
\]
and continuity of \(R\) give \(G_1(q_n)\to G_1(p)\).
The step-trial bounds therefore imply
\[
 \limsup_{N\to\infty}\E P_N\le S(p)+G_1(p).
\]
Trials with infinite
\(S(p)\) impose no further restriction, and taking the infimum
proves the proposition.
\end{proof}

The field supremum remains inside the trial infimum: the contact
argument fixes \(q\) before introducing its finite-dimensional
minimum. Section~\ref{sec:cavity} supplies the matching lower
bound, and Section~\ref{sec:general} passes from finite alphabets
to the spectral laws of Theorem~\ref{thm:main}.
